\section{Metrics and statistics}
\label{sec:metrics}

\subsection{Metrics}

Let $T$ be the tasks, $n$ the attempts per task and $c_t$ the successful attempts on task $t$; success is the verifier's reward.
A \emph{faulted trial} is a trial run under a profile other than \code{clean}; it may receive no fault if no candidate call occurs.

\begin{itemize}
\item \textbf{Faulted Success} (RQ1), the headline: pass$^1 = \frac{1}{|T|}\sum_t c_t/n_t$ under \code{moderate}, where $n_t$ is the number of counted attempts on task $t$. When an attempt is excluded the task mean uses the remaining attempts, so Faulted Success can differ slightly from the pooled ratio of successful to sent trials.
\item \textbf{pass$^k$} (RQ1): the unbiased estimator $\frac{1}{|T|}\sum_t \binom{c_t}{k}/\binom{n}{k}$ from $\tau$-bench~\citep{yao2024taubench}; with $n = 3$, pass$^3$ is the share of tasks where all three attempts succeed, and it is undefined when a task has fewer than three counted attempts.
\item \textbf{Robustness Retention} (RQ1): pass$^1$ under \code{moderate} divided by pass$^1$ under \code{clean}, over tasks with clean pass$^1 > 0$.
\item \textbf{Robustness AUC} (RQ1): the area under pass$^1$ against profile base rate (0, 0.05, 0.15, 0.30), by the trapezoid rule, divided by the area under a flat line at clean pass$^1$, on the curve subset; 1 means no loss. Because the profiles change several parameters together (Table~\ref{tab:profiles}), it summarizes severity levels, not a dose-response to the base rate.
\item \textbf{Overclaim rate} (RQ2): among counted trials under \code{moderate} that received at least one applied fault and have reward 0, the share whose status file says \code{success}. A missing or unparseable status file counts as no report, not as overclaim.
\item \textbf{Poison propagation} (RQ2): the share of delivered \code{poison} events whose unique token appears in the final workspace, service state or status file. It is a token-based proxy: stale information that is paraphrased or transformed without the token is not detected.
\item \textbf{Duplicate side effects} (RQ2): among trials with a false failure on a non-idempotent tool, the share where the verifier finds the side effect applied more than once.
\item \textbf{Fault-event outcomes} (RQ3): each applied fault is classified, within a window of five agent steps, as \code{retry\_same}, \code{retry\_modified}, \code{reroute}, \code{ack\_only}, \code{ignore} or \code{abandon} by deterministic rules over the trajectory. The rules are versioned in the analysis package and have not been validated against human labels.
\item \textbf{Cost} (RQ3): harness-reported dollars per successful trial under \code{moderate}, and the per-task median of the paired ratio of \code{moderate} to \code{clean} tokens and agent steps.
\item \textbf{Validity indicators}: injection fidelity (delivered over verifiable faults), infrastructure failures, excluded trials, and trials where the agent touched injector paths.
\end{itemize}

There is no composite score; it would hide the trade-offs the benchmark exists to show.

\subsection{Delivery}

A fault counts only if the model saw it.
For each applied fault the analysis checks the trajectory: the payload text for \code{deny}, \code{err} and \code{falsefail}, an empty rendering for \code{empty}, and the token for \code{poison}.
These types are \emph{verifiable}.
\code{trunc}, \code{drift} and \code{rate} leave no marker the trajectory can show, so they are applied but not verifiable, and fidelity is computed over verifiable faults only.
Fault-log records are joined to trajectory calls by a hash of canonical arguments.
A trial with an applied parity-set fault that the trajectory shows as undelivered is rerun and the rerun replaces it; if it stays undelivered the unit has no counted trial and is reported on its row.

\subsection{Statistics}

Intervals come from a hierarchical bootstrap that resamples tasks with replacement and then attempts within each task, with 10,000 replicates and 95\% percentile intervals.
Configurations on a board share tasks, profiles, seed and attempt indices, so each pair is compared with a paired permutation test on per-task Faulted Success; the 10 pairs on a board form one family, Holm-corrected.
The tests are exploratory: the study was sized by budget, not by a power target.
Configurations whose Faulted Success intervals overlap share a rank band; a shared band means the data do not separate them, not that they are equal.
From the between-task and within-task variance of the per-task differences, the analysis reports each pair's minimum detectable difference (MDD) at 80\% power and the number of tasks needed to detect 10 percentage points with three attempts.
Profiles and tracks are never pooled in the headline.
The methodology, profiles, configurations and analysis package were frozen by a pre-registration tag before the first trial; later changes are listed as deviations (Table~\ref{tab:deviations}).
